\documentclass[11pt]{article}
\usepackage[T1]{fontenc}
\usepackage{lmodern}
\usepackage{microtype}
\usepackage[margin=1in]{geometry}
\usepackage{amsmath,amssymb,amsthm,mathtools}
\usepackage{booktabs,longtable,array,enumitem}
\usepackage{xcolor,listings}
\usepackage[unicode,colorlinks=true,linkcolor=blue!45!black,urlcolor=blue!45!black,citecolor=blue!45!black]{hyperref}
\usepackage{fancyhdr}
\setlength{\headheight}{14pt}
\pagestyle{fancy}
\fancyhf{}
\fancyhead[L]{\small Variable-only equality for iPRA}
\fancyhead[R]{\small Research report}
\fancyfoot[C]{\thepage}
\setlength{\parskip}{0.35em}
\setlength{\parindent}{0pt}
\setlist{itemsep=0.2em,topsep=0.3em}
\newtheorem{theorem}{Theorem}[section]
\newtheorem{lemma}[theorem]{Lemma}
\newtheorem{corollary}[theorem]{Corollary}
\newtheorem{proposition}[theorem]{Proposition}
\theoremstyle{definition}
\newtheorem{definition}[theorem]{Definition}
\theoremstyle{remark}
\newtheorem{remark}[theorem]{Remark}
\DeclareMathOperator{\nf}{nf}
\newcommand{\ipra}{\mathsf{iPRA}}
\newcommand{\CV}{\mathcal C_V}
\newcommand{\CZ}{\mathcal C_0}
\newcommand{\CT}{\mathcal C_T}
\newcommand{\code}[1]{\texttt{\detokenize{#1}}}
\lstset{basicstyle=\ttfamily\scriptsize,columns=fullflexible,keepspaces=true,
  breaklines=true,showstringspaces=false,frame=single,rulecolor=\color{black!20},
  commentstyle=\color{black!50},keywordstyle=\color{blue!50!black},
  stringstyle=\color{teal!60!black},language=Python}
\title{\textbf{Variable-only equality substitution for iPRA}\\
  \large A five-rule calculus and a smaller complete normalized checker}
\author{Prepared for Vladimir Reshetnikov}
\date{November 11, 2026}
\begin{document}
\maketitle
\begin{abstract}
We further restrict the equality rule in a computation-normal presentation of
full first-order intuitionistic primitive recursive arithmetic. The rule
needs only substitution for a free numerical variable; it accepts no
term template and performs no arbitrary-term occurrence search. A short
construction packages two equations in one existential witness and recovers
every unrestricted crossed-transport inference. A proof-translation invariant
then shows that even the newest numerical parameter suffices.
The resulting Eq/$D$ calculus has five inference patterns. Its native
Eq/$Q$/$D$ companion has six and gives the smaller measured implementation.
Both are equiconsistent and bi-interpretable with iPRA, with whole numerical
domains, actual equality, and the identity numerical composites.
The new theorem concerns the equality inference. Computation normalization,
the paired binder, and the arithmetic interpretation bridge are inherited
and explicitly accounted for. Measurements compare complete standalone
checkers; no global minimality or new proof-assistant formalization is claimed.
\end{abstract}
\setcounter{tocdepth}{1}
\tableofcontents
\newpage

\section{Result, comparison boundary, and what is new}

The mathematical reduction is
\[
\frac{s=t\qquad h(s)=r}{r=h(t)}
\quad\rightsquigarrow\quad
\frac{x=t\qquad u=v}{v=\nf(u[t/x])},
\qquad x\text{ a numerical variable}.
\]
The right-hand rule is ordinary capture-free variable substitution followed
by the same computation normalization already required for binder instances.
The left-hand rule permits an arbitrary source term and a separately supplied
term template. We prove that both calculi derive exactly the same sequents.
The stronger refinement fixes $x$ to the newest numerical parameter of the
current proof scope.

The comparison is pinned to ProveIt commit
\[
\texttt{5d85023b0688e03b75f0ec3a04bb43ef0bb51705}.
\]
The project README at that commit describes the native scalar core, its
formalized interpretation into iPRA, and the existing reductions
\cite{repo}. We also inspected the five archives then in
\code{ConciseIPRA/docs/incoming}. In particular, deleting an Eval case,
removing $Q$ in favor of clauses, adopting full computation normal forms,
and using deterministic replacement of arbitrary source terms were already
represented by the prior packages. They are not attributed to this report.

The distinction from \code{template_free} is precise. That package replaces
all occurrences of an arbitrary term and recovers selected-context congruence
through a large raw computation marker. Full normalization erases such
markers. Here a numerical eigenvariable supplies the name instead.
The restriction to variable sources removes the need for arbitrary-term
matching altogether. The main proof goes further than the previously
proved exclusion of the literal-zero source.

The contribution is local to this pinned comparison. We do not claim
that the naming argument or its general proof-theoretic idea is historically
unprecedented. The statement that this exact restricted rule is sufficient
is proved below and is the claim on which the deliverable rests.

\paragraph{How to read the result.}
The five-rule presentation gives the cleanest formula language: only equality
and $D$. The six-rule companion retains $Q$, which makes its complete
goal-directed implementation smaller. Rule count and checker size therefore
favor different presentations. The implementation section reports the
coordinates separately instead of combining the best coordinates of
different programs.

\paragraph{Scope of the equivalence.}
Here \emph{iPRA} means full first-order intuitionistic logic over the primitive
recursive functions, with their defining equations, the usual successor
axioms, and induction for quantifier-free formulas. This is not merely
an equational fragment or a statement about standard natural numbers.
All interpretation domains used below are the entire numerical sort.
The conclusions include both interpretation composites, not just mutual
interpretability or equiconsistency.

